\section{La cardinalit\`a del numerabile}

\begin{definition}[Numerabilit\`a]
	Diciamo che $A$ \`e \vocab{al pi\`u numerabile} se $|A| \leq |\omega|$ ed \`e \vocab{numerabile} se $|A| = |\omega|$.
	Il simbolo $\aleph_0$ - aleph con zero - \`e semplicemente un'abbreviazione per $|\omega|$ (per cui $|A| \leq \aleph_0$ si pu\`o leggere ``$A$ \`e al pi\`u numerabile'' e $|A| = \aleph_0$ si 
	pu\`o leggere ``$A$ \`e numerabile'').
\end{definition}

\begin{remark}
	In altri termini, dire che $A$ \`e al pi\`u numerabile significa dire che c'\`e una funzione iniettiva $A \hookrightarrow \omega$. Dire che \`e numerabile significa dire che c'\`e una bigezione con $\omega$.
\end{remark}

\begin{proposition}[Dicotomia della al pi\`u numerabilit\`a]
	Se $A$ \`e al pi\`u numerabile, allora o $A$ \`e finito o $A$ \`e numerabile.
\end{proposition}

\textcolor{MidnightBlue}{Ossia $|A| < \aleph_0$ se e solo se $A$ \`e finito.}\\
Potremmo dimostrare la proposizione direttamente, ma ci conviene, invece, passare attraverso alcune considerazioni che saranno utili in seguito.\\
In generale, per costruire una bigezione fra due insiemi $A$ e $B$ - ossia per dimostrare $|A| = |B|$ - occorre appoggiarsi a qualche struttura definita sugli insiemi $A$
e $B$. Per esempio, una funzione successore. In questo corso, giocheranno un ruolo importante, in questa direzione, le relazioni d'ordine, e, in particolare - l'idea \`e di Cantor - i 
\vocab{buoni ordini}. Ricordiamo la definizione.

\begin{definition}[Buon ordinamento]
	Un insieme totalmente ordinato $(S,<)$ si dice \vocab{bene ordinato} se ogni suo sottoinsieme non vuoto ha un minimo.
	\[ \forall A \subseteq S \; A \ne \emptyset \rightarrow \exists m \in A \; \forall a \in A \; m \leq a
		\]
\end{definition}

Il trucco \`e che un isomorfismo di ordini \`e, in particolare, una bigezione, e spesso, per costruire bigezioni, costruiamo isomorfismi di ordini.

\begin{definition}[Isomorfismo]
	Due insiemi (parzialmente\footnote{Dove parziale indica l'assenza della propriet\`a di totalit\`a nella definizione di relazione d'ordine.}) ordinati
	$(A,<_A)$ e $(B,<_B)$ sono \vocab{isomorfi}, in simboli $(A,<_A) \sim (B,<_B)$ se esiste una bigezione $f : A \rightarrow B$ tale che:
	\[ \forall x,y \in A \; x <_A y \; \textcolor{purple}{\leftrightarrow} \; f(x) <_B f(y)
		\]
	cio\`e se esiste una bigezione che rispetta le relazioni d'ordine.
\end{definition}

\begin{remark}[Isomorfismi di ordini totali]
	Due insiemi \textcolor{red}{totalmente} ordinati $(A,<_A)$ e $(B,<_B)$ sono isomorfismi se e solo se esiste una funzione $f : A \rightarrow B$ \textbf{surgettiva} e \vocab{strettamente crescente} - cio\`e tale che:
	\[ \forall x,y \in A \; x <_A y \; \textcolor{purple}{\rightarrow} \; f(x) <_B f(y)
		\]
	(quando entrambi gli ordinamenti sono totali ci basta una sola freccia per avere in automatico l'altra, dunque diciamo che basta solo la stretta crescenza - che tecnicamente \`e solo la freccia sopra e non anche quella da destra verso sinistra).
\end{remark}

\begin{exercise}
	Dimostrare la proposizione enunciata sopra.
\end{exercise}

\begin{remark}[Ogni insieme finito \`e isomorfo alla sua cardinalit\`a]
	Sia $(A,<_A)$ totalmente ordinato con $|A| = n \in \omega$. Allora $(A,<_A) \sim (n,<)$, dove $<$ denota il buon ordinamento indotto da $\omega$.
\end{remark}

\begin{proof}
	Procediamo per induzione su $n$.
	\begin{itemize}
		\item[$\boxed{\text{caso $n = 0$}}$] $A = \emptyset$, quindi $(A,<_A) \sim (\emptyset,\emptyset)$ tramite la funzione vuota.
		\item[$\boxed{\text{caso $n = m + 1$}}$] Se $m = 0$, allora $A = \{a\}$ e $(A,<_A) \sim (1,<)$. Assumiamo quindi $m > 0$, e supponiamo che ogni insieme finito abbia un massimo.
		Sia $N$ il massimo di $A$, allora  $|A'| = m$, con $A' := A \setminus\{N\}$ (siamo nel caso $m > 0$, quindi \`e tutto legittimo), per cui per ipotesi induttiva esiste $f : A' \to m$ isomorfismo e quindi possiamo definire:
		\[ f' : A \to m + 1 : x \mapsto \begin{cases}
			f(x) &\text{se $x \in A'$} \\
			m &\text{se $x = N$}
		\end{cases}
			\]
		\`e banale che questa funzione sia ben definita e strettamente crescente.\\
		Ci resta da verificare l'assunzione iniziale, ovvero che \textcolor{purple}{ogni insieme finito e totalmente ordinato ha un massimo}.
		Procedendo per induzione, il caso $n = 0$ \`e banale, supponendo ora che $|A| = n$ ha massimo, vediamo che anche $|B| = n + 1$ ha massimo.
		Fissata $g$ bigezione tra $n + 1$ e $B$, abbiamo che $g_{|n}$ \`e una bigezione tra $n$ e $B \setminus\{g(m)\}$, pertanto, per ipotesi induttiva, quest'ultimo ammette massimo $M$.
		Ora per la totalit\`a dell'ordinamento su $B$ vale esattamente una tra $M > f(m)$ e $M < f(m)$, per cui [si verifica che] nel primo caso $M$ \`e il massimo di $B$, nel secondo $f(m)$ \`e il massimo di $B$,
		in ogni caso $B$ ha sempre un massimo, come voluto.
	\end{itemize}
\end{proof}

\begin{remark}[Ogni ordine totale finito un ha massimo e minimo]
	Questa proposizione ci dice che ogni ordine totale finito \`e isomorfo ad un buon ordine, $n \in \omega$, dunque ogni ordine totale finito ammette sia minimo - perch\'e $\omega$ \`e ben ordinato -,
	sia massimo - perch\'e l'abbiamo dimostrato nella dimostrazione precedente in generale, o anche come conseguenza della caratterizzazione di $(\omega,<)$ che stiamo per vedere.
\end{remark}

Possiamo caratterizzare $\omega$ in termini delle propriet\`a del suo ordinamento naturale, grazie alle propriet\`a seguenti.

\begin{proposition}[Propriet\`a di $(\omega,<)$ come ordine totale]
	Dato $(\omega,<)$ ordine totale allora valgono le seguenti:
	\begin{enumerate}[(1)]
		\item $(\omega,<)$ \`e un buon ordine.
		\item $(\omega,<)$ \`e \vocab{illimitato} - ossia $\forall x \in \omega \; \exists y \in \omega \; x < y$.
		\item Ogni $A \subseteq \omega$ superiormente limitato e non vuoto ha un massimo, ossia:
		\[ \forall A \subseteq \omega (A \ne \emptyset \land (\exists L \in \omega \; \forall x \in A \; x \leq L)) \rightarrow (\exists M \in A \; \forall x \in A \; x \leq M)
			\]
	\end{enumerate}
\end{proposition}

\begin{proof}
	Abbiamo che (1) \`e il principio del minimo che abbiamo gi\`a dimostrato su $\omega$, per (2) basta prendere $y = s(x)$ (e $x \in s(x) \implies x < y$). Per (3) se $A$ \`e superiormente limitato da $L \in \omega$, allora $A \subseteq s(L)$, quindi $A$ \`e finito - perch\'e sottoinsieme di un insieme finito.
	Siccome $A$ \`e finito ha un massimo per quanto osservato in precedenza, oppure perch\'e l'ordinamento su $A$ definito da:
	\[ x \prec y \Mydef x > y \,\footnote{Noto anche come \vocab{ordinamento duale} o \vocab{duale} dell'ordinamento $<$.}
		\]
	\`e totale, per la totalit\`a di $<$, ed \`e anche buono perch\'e ogni sottoinsieme non vuoto di $A$ \`e finito, totalmente ordinato, ed ha ha massimo per $<$ per l'osservazione precedente, che \`e dunque minimo per $\prec$.
	Per cui $(A,\prec)$ \`e un buon ordinamento, quindi ha un minimo per la relazione $\prec$, che corrisponde ad un massimo per $<$.\footnote{In ogni caso facciamo sempre affidamento al fatto che ogni insieme finito abbia massimo.}
\end{proof}

\begin{proposition}[Caratterizzazione di $(\omega,<)$ come ordine totale]
	Sia $(A,\prec)$, con $A \ne \emptyset$, un ordinamento:
	\begin{enumerate}
		\item buono
		\item illimitato
		\item tale che ogni sottoinsieme superiormente limitato e non vuoto di $A$ ha un massimo secondo $\prec$
	\end{enumerate}
	allora $(A,\prec) \sim (\omega,<)$.\footnote{Questa proposizione completa la caratterizzazione di $(\omega,<)$ come ordine totale.}
\end{proposition}

Dimostriamo prima un facile lemma.

\begin{lemma}[Stretta crescenza col successore $\implies$ stretta crescenza]
	Sia $(A,\prec)$ un ordine (non necessariamente totale), e sia $f : \omega \rightarrow A$ tale che:
	\[ \forall n \in \omega \; f(n) \prec f(s(n)) \, \footnote{Typo di Mamino nelle dispense su $\prec$ e $<$.}
		\]
	allora $f$ \`e strettamente crescente, cio\`e $\forall m,n \in \omega \; m < n \rightarrow f(m) \prec f(n)$.
\end{lemma}

\begin{proof}
	Consideriamo per assurdo $m < n$ tali che $f(m) \not \prec f(n)$, con $n$ minimo. Siccome $0 \leq m < n$, esiste $n'$ tale che $n = n' + 1$, con $m \geq n'$.
	Nel primo caso $f(m) \overset{\text{Hp.}}{\prec} f(s(m)) = f(n)\;\textcolor{red}\lightning$, e analogamente nel secondo:
	\[ f(m) \overset{\text{$n$ minimo}}{\preceq} f(n') \overset{\text{Hp.}}{\prec} f(n) \;\textcolor{red}\lightning
		\]
	abbiamo quindi in tutti i casi un assurdo.
\end{proof}

Possiamo ora dimostrare la proposizione.\\
\textcolor{MidnightBlue}{L'idea \`e che le propriet\`a elencate per ipotesi siano il minimo indispensabile che un ordinamento totale deve avere affinch\'e si possa sempre costruire un isomorfismo tra $\omega$ e un tale ordinamento totale.}

\begin{proof}
	Costruiamo per ricorsione prima forma un isomorfismo $f$ da $(\omega, <)$ a $(A,\prec)$ nella maniera seguente:
	\[ f(0) = \min_\prec A \qquad f(s(n)) = \min_\prec\{a \in A | f(n) \prec a\} \,\footnote{Cio\`e la funzione prende ogni volta l'elemento pi\`u piccolo non ancora nell'immagine, come vedremo questo \`e l'unico isomorfismo sensato - e possibile - tra buoni ordinamenti.}
		\]
	dove $\min_\prec$ denota il minimo secondo la relazione d'ordine (buona) $\prec$ di $A$. Verifichiamo che $f : \omega \to A$ \`e ben definita, strettamente crescente e surgettiva.
	\begin{itemize}
		\item[$\diamondsuit$] \underline{$f$ ben definita}: il teorema di ricorsione numerabile ci garantisce esistenza ed unicit\`a di $f$ a patto che la $h : A \to A : x \mapsto \min_\prec\{x \in A | x \prec a\}$ sia ben definita, in questo caso i minimi sono ben definiti perch\'e per ipotesi (1) $A$ \`e ben ordinato, nel caso di $f(0)$ naturalmente assumiamo che $A \ne \emptyset$,
		mentre per il successore osserviamo che $\{a \in A | f(n) \prec a\}$ \`e sempre non vuoto, per ogni $n \in \omega$, perch\'e per ipotesi (2) $A$ \`e superiormente illimitato - se l'insieme fosse vuoto allora $A$ sarebbe superiormente limitato da un suo elemento.
		\item[$\diamondsuit$] \underline{$f$ strettamente crescente}: per costruzione $f(n) \prec f(s(n))$ per ogni $n \in \omega$, per cui $f$ soddisfa le ipotesi del lemma precedente ed \`e strettamente crescente (quindi in particolare iniettiva).
		\item[$\diamondsuit$] \underline{$f$ surgettiva}: dato $y \in A$ cerchiamo $x \in \omega$ tale che $f(x) = y$. Scegliamo $x := \min_{\prec}\{n \in \omega | y \prec f(n)\}$, possiamo farlo perch\'e $\omega$ \`e bene ordinato e se per assurdo $y \succ f(n)$ per ogni $n \in \omega$,
		allora $f[\omega]$ sarebbe un sottoinsieme di $A$ superiormente limitato e senza massimo - basta prendere $f$ del successore -, contraddicendo (3). Abbiamo quindi $y \prec f(x)$, osserviamo che vale necessariamente anche la disuguaglianza opposta ed abbiamo concluso, distinguiamo due casi.
		\begin{itemize}
			\item[$\boxed{x = 0}$] In tal caso $f(x) = \min_\prec A$, quindi necessariamente $f(x) \preceq y$.
			\item[$\boxed{x = x' + 1}$] In questo caso, per la minimalit\`a di $x$, si ha che $f(x') \preceq y$, naturalmente se $f(x') = y$ abbiamo comunque concluso\footnote{Typo (?), Mamino questo caso l'ha saltato.}, per cui ci rimane il caso $f(x') \prec y$.
			Avendo definito $f(x) = f(s(x'))$ come il minimo in $A$ pi\`u grande di $f(x')$, segue che $f(x) \preceq y$ - $y$ \`e nell'insieme in cui prendiamo tale minimo.
		\end{itemize}
	\end{itemize}
\end{proof}

Tornando alla proposizione iniziale.

\begin{proposition}[Dicotomia della al pi\`u numerabilit\`a]
	Se $A$ \`e al pi\`u numerabile, allora o $A$ \`e finito o $A$ \`e numerabile.
\end{proposition}

\begin{proof}
	Per ipotesi esiste $f : A \rightarrow \omega$ iniettiva, per cui abbiamo $|A| = |f[A]|$, e siccome $f[A] \subseteq \omega$,
	ci basta dimostrare che dato $B \subseteq \omega$, o $B$ \`e finito o \`e numerabile.\\
	Dato un sottoinsieme $B \subseteq \omega$ o \`e finito o \`e infinito, nel primo caso abbiamo gi\`a concluso, nel secondo mostriamo che $(B,<_{|B}) \sim (\omega,<)$, e di conseguenza $|B| = \aleph_0$.\\
	Naturalmente $(B,<_{|B})$ continua ad essere un buon ordinamento e qualsiasi sottoinsieme superiormente limitato continua ad avere massimo grazie alle propriet\`a di $\omega$\footnote{Typo di Mamino che scambia 2 e 3 nella dimostrazione.} - i sottoinsiemi di $B$ sono in particolare sottoinsiemi di $\omega$ -,
	per dimostrare che vale (2) dobbiamo invece verificare che $B$ non ha massimo elemento. Se per assurdo $M:=\max B$ allora $B \subseteq s(M)$, ma $s(M)$ \`e finito, per cui anche $B$ lo sarebbe, che \`e contro l'ipotesi.
\end{proof}

\begin{exercise}[Disuguaglianza con la surgettivit\`a senza AC]
	\label{disuguaglianze_senza_AC2}
	Dimostra che se $|A| \leq \aleph_0$ e $f : A \twoheadrightarrow B$ \`e surgettiva, allora $|B| \leq \aleph_0$.
\end{exercise}

\begin{soln}
	Mostriamo che sotto queste ipotesi esiste $h : B \hookrightarrow \omega$ - iniettiva. Sia $g : A \hookrightarrow \omega$ e poniamo:
	\[ h(b) = \min_<(g[f^{-1}(b)])
		\]
	$f^{-1}(b) \ne \emptyset$ per ogni $b \in B$ poich\'e $f$ \`e surgettiva, per cui il minimo \`e ben definito, e quindi lo \`e $h$.
	Vediamo l'iniettivit\`a, se $h(b) = h(b')$, allora i minimi, siano $g(a)$ e $g(a')$, sono uguali, ma $a$ e $a'$ sono elementi nelle controimmagini rispettivamente di $b$ e $b'$, cio\`e tali che $f(a) = b$ e $f(a') = b'$.
	Sappiamo quindi per ipotesi che $g(a) = g(a')$ e per l'iniettivit\`a di $g$ segue $a = a'$, da cui $f(a) = f(a')$, da cui $b = f(a) = f(a') = b'$.
\end{soln}

\subsection{Insiemi numerabili in pratica}
Sapere che, se $|A| \leq \aleph_0$, allora o $A$ \`e finito o \`e numerabile, ci fornisce lo strumento essendo per dimostrare la numerabilit\`a di molti insiemi concreti. Spesso, infatti,
\`e facile dimostrare che un insieme infinito \`e tale. Rimane poi da gestire un discorso di disuguaglianze per dire che esso \`e al pi\`u numerabile.\\
Cominciamo quindi con qualche considerazione generale a proposito delle disuguaglianze fra cardinalit\`a.

\begin{remark}[Compatibilit\`a tra operazioni e ``ordinamentro'' fra cardinalit\`a]
	\label{compatibilit\`a_operazioni_cardinalit\`a}
	Dati gli insiemi $A,B,C$ con $|B| \leq |C|$ allora vale:
	\begin{align*}
		|A| + |B| \leq |A| + |C| \qquad & |A|^{|B|} \leq |A|^{|C|} \\
		|A| \cdot |B| \leq |A| \cdot |C| \qquad & |B|^{|A|} \leq |C|^{|A|}
	\end{align*}
\end{remark}

Vale a dire che le operazioni sulle cardinalit\`a sono monotone, nel senso delle disuguaglianze larghe.
\textcolor{red}{Attenzione per\`o che, in generale, NON sono strettamente monotone!}

\begin{proof}
	Detta $f : B \rightarrow C$ la funzione iniettiva che testimonia che $|B| \leq |C|$ e detto $B' = f[B]$ abbiamo che $|B| = |B'|$,
	quindi basta dimostrare le disuguaglianze asserite con $B'$ al posto di $B$\footnote{Oppure potevamo assumere WLOG che $B$ fosse proprio contenuto in $C$ e che la mappa
	fosse proprio $\id_B$, in ogni caso \`e solo una questione di nomi.}. Ora, poich\'e $B' \subseteq C$ otteniamo facilmente:
	\[ \begin{split}
		B' \subseteq C \overset{\text{ovvio}}{\implies}& (A \times \{0\}) \cup (B' \times \{1\}) \subseteq (A \times \{0\}) \cup (C \times \{1\}) \; \textcolor{red}{= A \sqcup B' \subseteq A \sqcup C}\\
					   \overset{\id_A \times \id_{B'}}{\implies}& |(A \times \{0\}) \cup (B' \times \{1\})| \leq |(A \times \{0\}) \cup (C \times \{1\})| \\
					   \overset{\text{def.}}{\iff}& |A| + |B'| \leq |A| + |C|
	\end{split}
		\]
	Le altre si ottengono allo stesso modo.
\end{proof}

\begin{remark}[Disuguaglianza di inclusione-esclusione]
	$|A \cup B| \leq |A| + |B|$.
\end{remark}

\begin{proof}
	Basta osservare che la seguente funzione \`e iniettiva:
	\[ f : A \cup B \to (A \times \{0\}) \cup (B \times \{1\}) : x \mapsto \begin{cases}
		(x,0) &\text{se $x \in A$} \\
		(x,1) &\text{altrimenti}
	\end{cases}
		\]
\end{proof}

Veniamo ora a calcolare le operazioni aritmetiche. Gi\`a sappiamo, per il \hyperref[cantor]{teorema di cantor}, 
che $2^{\aleph_0} > \aleph_0$, per cui mettere un $\aleph_0$ a esponente di qualunque cosa non sia uno 0 o un 1 conduce fuori dal numerabile.
Tutto il resto invece no.

\begin{proposition}[Operazioni aritmetiche con $\aleph_0$]
	$\aleph_0 + \aleph_0 = \aleph_0 \cdot \aleph_0 = \aleph_0^{n} = \aleph_0$, con $n \in \omega\setminus\{0\}$.
\end{proposition}

\begin{proof}
	Supponiamo di sapere gi\`a che $\aleph_0 \cdot \aleph_0 = \aleph_0$, allora possiamo formare la catena di disuguaglianze:
	\[ \aleph_0 \overset{\text{op. card.}}{=} \aleph_0 + 0 \overset{\text{oss. sopra}}{\leq} \aleph_0 + \aleph_0 \overset{\text{op. card.}}{=} \aleph_0 \cdot 2 \overset{\text{oss. sopra}}{\leq} \aleph_0 \cdot \aleph_0 \overset{\text{ipotesi}}{=} \aleph_0
		\]
	Da cui per il teorema di \hyperref[CB]{Cantor-Bernstein}:
	\[ \aleph_0 + \aleph_0 = \aleph_0 \cdot \aleph_0 = \aleph_0
		\]
	Ora \`e facile vedere per induzione che $n \in \omega \setminus\{0\} \rightarrow \aleph_0^n = \aleph_0 $, infatti $\aleph_0^1 = \aleph_0$ [e $\aleph_0^2 = \aleph_0 \cdot \aleph_0 = \aleph_0$], quindi $\aleph_0^{n+1} = \aleph_0^n \cdot \aleph_0 \overset{\text{Hp. indutt.}}{=} \aleph_0 \cdot \aleph_ 0 = \aleph_0$.
\end{proof}

Per concludere la dimostrazione precedente, resta da dimostrare il lemma seguente.

\subsection{Prodotto di numerabili \`e numerabile}

\begin{lemma}[$\aleph_0 \cdot \aleph_0 = \aleph_0$]
	$\aleph_0 \cdot \aleph_0 = \aleph_0$, ossia esiste una bigezione fra $\omega \times \omega$ e $\omega$.
\end{lemma}

\begin{wrapfigure}[15]{r}{2.82cm}
	\includegraphics[width=2.82cm]{immagini/aleph2.png}
\end{wrapfigure}
Ci sono diverse vie per illustrare questo risultato. Per esempio, possiamo rappresentare le coppie $(x,y) \in \omega \times \omega$ sotto la specie di una griglia a maglie quadrate.
Poi disegnare un percorso che pare visitare tutte le maglie della griglia, con sufficiente apparenza di regolarit\`a, possibilmente, da convincere il lettore che vi debba essere un metodo.
Infine numeriamo le maglie secondo l'ordine in cui sono visitate dal percorso. Avremo cos\`i numerato tutte le coppie di numeri naturali del disegno.\\
Altrimenti, \`e possiamo esibire delle bigezioni esplicite, per esempio:
\[ f(x,y) = 2^x \cdot (2y + 1) - 1 \qquad g(x,y) = \frac{(x+y)^2 + 3x + y}{2}
	\]
\`E possibile anche scrivere i due numeri della coppia in base 10 a cifre alternate, tipo: $(\textcolor{blue}{64},\textcolor{red}{4096}) \mapsto \textcolor{red}{4}\textcolor{cyan}{0}\textcolor{red}{0}\textcolor{red}{9}\textcolor{cyan}{0}\textcolor{blue}{6}\textcolor{red}{6}\textcolor{red}{4}\textcolor{blue}{4}$.\\

\newpage
\hspace{-0.43cm}\emph{Dimostrazione.}\hspace{-0.20cm} Innanzitutto definiamo su $\omega \times \omega$ un ordinamento come segue:
	\begin{wrapfigure}[9]{l}{2.5cm}
		\includegraphics[width=2.5cm]{immagini/ordine_omega.png}
	\end{wrapfigure}
	\[ \begin{split}
		(a,b) \prec (a',b') &\Mydef \max(a,b) < \max(a',b') \\
							&\lor (\max(a,b) = \max(a',b') \land a < a') \\
							&\lor (\max(a,b) = \max(a',b') \land a = a' \land b < b')
	\end{split}
		\]
	ossia per confrontare $(a,b)$ con $(a',b')$, si confrontano prima $\max(a,b)$ e $\max(a',b')$; a parit\`a si confrontano $a$ ed $a'$; se queste coincidono, allora si confrontano $b$ e $b'$.
	\textcolor{MidnightBlue}{L'idea \`e che le coppie $\prec$ di una certa $(a,b)$ fissata sono tutte contenute nel quadrato $\{0,\ldots,\max(a,b)\} \times \{0,\ldots,\max(a,b)\}$, per cui sono in numero finito, e quindi si ha $(\omega \times \omega, \prec)$ isomorfo a $(\omega,<)$.}\\
	Formalmente, iniziamo col verificare che $\prec$ sia effettivamente un ordine stretto e totale.
	\begin{itemize}
		\item \underline{irriflessivit\`a}: immediata dal fatto che confrontare $(a,b)$ con se stessa da un OR con tre alternative tutte necessariamente false a causa dell'irriflessivit\`a di $<$.
		\item \underline{transitivit\`a}: data $(a,b) \prec (a',b) \prec (a'',b'')$ vogliamo ottenere che $(a,b) \prec (a'',b'')$. Dalle disuguaglianze precedenti segue $\max(a,b) \leq \max(a',b') \leq \max(a'',b'')$, se una di queste disuguaglianze \`e stretta allora $(a,b) \prec (a'',b'')$ e si conclude,
		altrimenti $\max(a,b) = \max(a',b') = \max(a'',b'')$, e segue dalla definizione di $\prec$ che $a \leq a' \leq a''$. Nuovamente, se una delle disuguaglianze \`e strette abbiamo concluso, altrimenti per definizione abbiamo $b \leq b' \leq b''$, e necessariamente una delle disuguaglianze deve essere stretta,
		altrimenti all'inizio avremmo avuto un'uguaglianza, che \`e assurdo, per cui si conclude necessariamente.
		\item \underline{totalit\`a}: basta osservare che se per $(a,b)$ ed $(a',b')$ non vale nessuna delle due fra $\prec$ e $\succ$, allora necessariamente $a = a'$ e $b = b'$, che rende anche in questo caso vera la definizione di ordinamento totale.
	\end{itemize}
	Ora vogliamo dire che $(\omega \times \omega, \prec) \sim (\omega, <)$, in tal modo, avendo un'isomorfismo di ordini, avremmo in particolare una bigezione tra $\omega$ e $\omega \times \omega$.
	Partiamo dall'osservazione che se $(a,b) \in \omega \times \omega$ allora possiamo definire:
	\[ (\omega \times \omega)_{(a,b)} \Mydef \{(x,y) \in \omega \times \omega | (x,y) \prec (a,b)\}
		\]
	detto il ``\vocab{segmento iniziale} determinato da $(a,b)$ su $(\omega \times \omega,\prec)$''. Tale segmento iniziale \`e finito, infatti $(\omega \times \omega)_{(a,b)} \subseteq s(\max(a,b)) \times s(\max(a,b))$, dove il RHS \`e un insieme finito e quindi tutti i suoi sottoinsiemi sono finiti.\\
	Ci serve dire: \textcolor{red}{1.} $(\omega \times \omega, \prec)$ \`e bene ordinato \textcolor{red}{2.} $(\omega \times \omega, \prec)$ \`e illimitato \textcolor{red}{3.} ogni sottoinsieme non vuoto e superiormente limitato di $\omega \times \omega$ ha un massimo.
	\begin{enumerate}[1.]
		\item Dato $A \subseteq \omega \times \omega$ con $A \ne \emptyset$, possiamo fissare $a \in A$ e considerare $(\omega \times \omega)_a \cap A$. Se tale intersezione \`e vuota allora si vede banalmente che $a$ \`e il minimo di $A$. Se l'intersezione \`e non vuota, allora \`e un sottoinsieme di un insieme finito e totalmente ordinato, dunque ha minimo $m \in A$ e tale minimo 
		lo \`e proprio per tutti gli elementi di $A$, infatti, preso $x \in A \setminus((\omega \times \omega)_a \cap A)$, si ha $m < a \leq x$.
		\item Dato $(a,b) \in \omega \times \omega$, $(a,b) \prec (s(a),s(b))$, dunque $\omega \times \omega$ \`e illimitato.
		\item Dato $A \subseteq \omega \times \omega$ non vuoto e superiormente limitato da $(a,b) \in \omega \times \omega$, abbiamo che $A \subseteq (\omega \times \omega)_{(a+1,b+1)}$, dunque \`e finito, e quindi ammette massimo perch\'e $\prec$ \`e totale e valgono le osservazioni fatta in precedenza.
	\end{enumerate}
$\hfill\square$

\subsection{Numeri interi e razionali}
Usando la proposizione appena dimostrata, potremmo dimostrare, per esempio, che $\ZZ$ e $\QQ$ sono numerabili, se non fosse che non abbiamo ancora definito questi oggetti. Allo scopo, ricordiamo che - \hyperref[3.73]{esercizio} - una relazione di equivalenza
induce un insieme di classi di equivalenza.

\begin{definition}[$\ZZ$]
	Definiamo $\ZZ$ come l'insieme delle classi di equivalenza su $\omega \times \omega$ indotte dalla relazione:
	\[ (a,b) \sim_\ZZ (a',b') \Mydef a + b' = b + a' \, \footnote{Moralmente: ``$(a,b) = a - b$''.}
		\]
\end{definition}

\begin{exercise}
	Dimostrare che $\sim_\ZZ$ \`e una relazione di equivalenza.
\end{exercise}

\begin{example}[Operazioni su $\ZZ$]
	Definiamo $+,-,\cdot$ su $\ZZ$ mediante:
	\begin{align*}
		[(a,b)]_\ZZ + [(a',b')]_\ZZ &\Mydef [(a+a',b+b')]_\ZZ \\
		-[(a,b)]_\ZZ &\Mydef [(b,a)]_\ZZ \\
		[(a,b)]_\ZZ \cdot [(a',b')]_\ZZ &\Mydef [(a\cdot a' + b \cdot b', a \cdot b' + a' \cdot b)]_{\ZZ}
	\end{align*}
	dimostra che $\ZZ$, con queste operazioni, \`e un anello commutativo con identit\`a: $1 \Mydef [(1,0)]_\ZZ$. 
\end{example}

\begin{definition}[$\QQ$]
	Definiamo $\QQ$ come l'insieme delle classi di equivalenza su $\ZZ \times (\omega\setminus\{0\})$ indotte dalla relazione:
	\[ (n,d) \sim_\QQ (n',d') \Mydef n \cdot d' = n' \cdot d \, \footnote{Moralmente: ``$(n,d) = \frac nd$''.}
		\]
\end{definition}

\begin{exercise}
	Dimostrare che $\sim_\QQ$ \`e una relazione di equivalenza.
\end{exercise}

\begin{exercise}[Operazioni su $\QQ$]
	Definisci $+,-,\cdot$ e $\square^{-1}$ su $\QQ$ nella maniera ragionevole e dimostra che $\QQ$ \`e un campo.
\end{exercise}

\begin{exercise}[Ordinamento su $\QQ$]
	Definisci la relazione $<$ su $\QQ \times \QQ$ dicendo che $q \in \QQ$ \`e positivo se $q = [(n,d)]_\QQ$, con $n,d \in \omega\setminus\{0\}$, e dicendo che $a < b$ se e solo se $b - a$ \`e positivo.
	Dimostra che questo \`e un ordine totale e \vocab{denso}, cio\`e:
	\[ \forall a,b \in \QQ \; a < b \rightarrow \exists c \in \QQ \; a< c <b \, \footnote{Typo di Mamino.}
		\]
\end{exercise}

\begin{note}
	Gli esercizi precedenti sono tediosi, ma non sono difficili. Nel resto del corso daremo per scontate le propriet\`a aritmetiche elementari di $\ZZ$ e $\QQ$.
	D'ora innanzi scriveremo:
	\[ a - b \Mydef [(a,b)]_{\ZZ} \qquad \frac{n}{d} \Mydef [(n,d)]_\QQ
		\]
\end{note}

Per dimostrare la numerabilit\`a di $\ZZ$ e $\QQ$, \`e comodo richiamare ancora un \hyperref[disuguaglianze_senza_AC2]{esercizio}, per\`o, questa volta, lo risolviamo\footnote{La soluzione riportata all'esercizio di riferimento \`e quella di Mamino.}.

\begin{corollary}[Caratterizzazione della al pi\`u numerabilit\`a]
	\label{disugcardnum}
	Un insieme $A \ne \emptyset$ \`e al pi\`u numerabile se e solo se esiste $f : \omega \rightarrow A$ surgettiva.\footnote{D'ora in avanti le disuguaglianze che danno al pi\`u numerabile ottenute mediante mappe surgettive saranno quindi legittime.}
\end{corollary}

\begin{proof}
	La freccia $\Longleftarrow$ deriva dall'esercizio citato prima. Per l'inverso, supponiamo $A$ al pi\`u numerabile e mostriamo che c'\`e sempre una mappa surgettiva tra $\omega$ ed $A$.
	Per ipotesi esiste $g : A \hookrightarrow \omega$ - iniettiva -, essendo $A \ne \emptyset$ possiamo fissare $a \in A$ e definire la seguente mappa:
	\[ f : \omega \to A : x \mapsto \begin{cases}
										g^{-1}(x) &\text{se $x \in \Imm(g)$} \\
										a &\text{altrimenti}
									\end{cases}
		\]
	tale mappa \`e naturalmente ben definita ed \`e surgettiva in quanto stiamo semplicemente estendendo la mappa bigettiva $g^{-1} : \omega\supseteq\Imm(g) \to A$.
\end{proof}

\begin{notation}[Successione - enumerazione]
	Con \vocab{successione} (numerabile) intendiamo semplicemente una funzione con dominio $\omega$, per cui:
	\[ \alpha = \{\alpha_i\}_{i \in \omega} \Mydef \alpha : \omega \to \ldots: i \mapsto \alpha_i
		\]
	cio\`e $\alpha$ \`e un ``elenco numerabile'' - non necessariamente esaustivo - degli elementi dell'insieme di arrivo. Per \vocab{enumerazione}
	intendiamo una successione (numerabile) surgettiva, tale per cui, detto $A$ l'insieme di arrivo, si ha $\Imm(\alpha) = A$, o - informalmente - $A = \{\alpha_i | i \in \omega\}$.
\end{notation}

\textcolor{MidnightBlue}{Il corollario sopra, quindi, non ci dice altro che $A \ne \emptyset$ \`e al pi\`u numerabile se e solo se ha almeno un'enumerazione.}

\begin{example}[L'insieme dei numeri interi \`e numerabile]
	$\ZZ$ \`e numerabile.
\end{example}

\begin{proof}
	La funzione $\omega \times \omega \to \ZZ : (a,b) \mapsto a - b$ \`e surgettiva per definizione (\`e la proiezione al quoziente di $\omega \times \omega$ modulo $\sim_\ZZ$, che sappiamo essere sempre surgettiva, in questo caso stiamo indicando le classi $[(a,b)]_\ZZ$ con $a - b$, ma sono sempre classi di equivalenza), e $\omega \times \omega$ \`e numerabile\footnote{$|\omega \times \omega| = \aleph_0 \cdot \aleph_0 = \aleph_0$.} dunque $|\ZZ|\leq \aleph_0$.\\
	D'altro canto, la funzione $\omega \rightarrow \ZZ : n \mapsto [(n,0)]_\ZZ$ \`e iniettiva, infatti $[(n,0)]_\ZZ = [(m,0)]_\ZZ \iff (n,0) \sim (m,0) \iff n = m$ (per definizione di $\sim_\ZZ$), dunque $\aleph_0 \leq |\ZZ|$, pertanto - per \hyperref[CB]{Cantor-Bernstein} - $|\ZZ| = \aleph_0$.
\end{proof}

\begin{example}[L'insieme dei numeri razionali \`e numerabile]
	$\QQ$ \`e numerabile.
\end{example}

\begin{proof}
	Come nell'esempio precedente, la proiezione al quoziente $\ZZ \times (\omega \setminus\{0\}) \rightarrow \QQ : (n,d) \mapsto \frac{n}{d}$ (dove la frazione \`e un'abbreviazione per la classe di equivalenza $[(n,d)]_\QQ$), \`e surgettiva per costruzione, inoltre $|\ZZ \times (\omega \setminus\{0\})| = |\ZZ| \cdot |\omega \setminus\{0\}| = \aleph_0 \cdot \aleph_0 = \aleph_0$, dunque vale il 
	\hyperref[disugcardnum]{corollario} sulla disuguaglianza tra cardinalit\`a, pertanto $\aleph_0 \geq |\QQ|$.\\
	Viceversa, la funzione $\omega \rightarrow \QQ : n \mapsto \frac{n}{1}$ \`e iniettiva, infatti $\frac n1 = \frac m1 \iff n \cdot 1 = m \cdot 1 \iff m = n$, dunque per definizione si ha $\aleph_0 \leq |\QQ|$. Da cui per \hyperref[CB]{Cantor-Bernstein} $|\QQ| = \aleph_0$.
\end{proof}

Adesso, ci piacerebbe poter dire che, se abbiamo un insieme $A$ al pi\`u numerabile, e tutti i suoi elementi sono, a loro volta, insiemi al pi\`u numerabili, allora
$\bigcup A$ \`e al pi\`u numerabile \textcolor{MidnightBlue}{- ovvero un'unione al pi\`u numerabile di insiemi al pi\`u numerabili \`e a sua volta al pi\`u numerabile}.\\
D'altro canto \`e ragionevole: se esiste una enumerazione $\{a_i\}_{i \in \omega}$ di $A$, e, per ogni $i \in \omega$, esiste una enumerazione di $a_i$, cio\`e esiste:
\[ \alpha_i : \omega \twoheadrightarrow a_i : j \mapsto \alpha_i(j) = \alpha_{i,j}
	\]
tale per cui $a_i = \Imm(\alpha_i) = \{\alpha_{i,j}\}_{j \in \omega}$, allora possiamo mandare surgettivamente - cio\`e enumerare - $\omega \times \omega$ in $\bigcup A$ mandando ogni coppia $(i,j)$ nel $j$-esimo elemento dell'$i$-esimo elemento di $A$, $(i,j) \mapsto \alpha_{i,j}$,
e, siccome $\omega \times \omega$ numerabile, allora per il \hyperref[disugcardnum]{corollario} $|A| \leq \aleph_0$.\\
L'\textcolor{red}{errore} \`e credere di poter fissare una enumerazione $\alpha_i$ di $a_i$ per ogni $i \in \omega$. Usando l'assioma della scelta potremo farlo, ma, per ora, non abbiamo modo, in generale, di procurarci la funzione che associa ogni naturale ad una enumerazione dell'elemento di $A$ da lui indicizzato, $i \mapsto \alpha_i$.
Possiamo per\`o assumere di averla, cos\`i si corregge il ragionamento impreciso di prima.

\begin{proposition}[$|A| \leq \aleph_0 \implies |\bigcup A| \leq \aleph_0$]
	Sia $A = \{a_i \in A| i \in \omega\}$ e \textcolor{red}{sia $\{\alpha_i\}_{i \in \omega}$ una successione di funzioni}\footnote{Che stiamo appunto assumendo di avere gi\`a, altrimenti serve AC.} tali per cui,
	per ogni $i \in \omega$, si ha che $\alpha_i : \omega \twoheadrightarrow a_i$ \`e una enumerazione di $a_i$, allora $\left|\bigcup A\right| \leq \aleph_0$.
\end{proposition}

\begin{proof}
	Basta osservare che la funzione:
	\[ f : \omega \times \omega \to \bigcup A : (i,j) \mapsto \alpha_i(j) = \alpha_{i,j}
		\]
	\`e surgettiva e vale quindi il solito \hyperref[disugcardnum]{corollario}. Infatti, dato $a \in \bigcup A$,
	poich\'e $A = \{a_i\}_{i \in \omega}$, allora $a \in a_i$, per qualche $i \in \omega$, e poich\'e $a_i = \Imm(\alpha_i)$ - qui stiamo usando che gli elementi di $A$ sono a loro volta al pi\`u numerabili -,
 	allora esiste $j \in \omega$ tale per cui $\alpha_i(j) = a$.
\end{proof}

\pagebreak

\begin{notation}
	Data una funzione $f : I \rightarrow S$ definiamo:
	\[ \bigcup_{i \in I} f(i) \Mydef \bigcup f[I]
		\]
	Cos\`i, per esempio, se $A = \{a_i | i \in \omega\}$:
	\[ \bigcup A = \bigcup_{i \in \omega} a_i \;\textcolor{MidnightBlue}{= \{x | \exists i \in \omega \; x \in a_i\}}
		\]
\end{notation}

\begin{definition}
	[Parti finite]
	Definiamo le \vocab{parti finite} di un insieme $A$ come:
	\[ \psf(A) \Mydef \{X \in \ps(A) : |X| < \aleph_0\}
		\]
\end{definition}

\begin{proposition}[Insieme al pi\`u numerabile $\implies$ parti finite al pi\`u numerabile]
	$|A| \leq \aleph_0 \rightarrow |\psf(A)| \leq \aleph_0$.
\end{proposition}

\begin{proof}
	Il caso $A = \emptyset$ \`e immediato. Assumiamo $A \ne \emptyset$, sia:
	\[ \ps^{\leq n} = \{X \in \ps(A) : |X| \leq n\}
		\]
	siccome $\psf(A) = \bigcup_{n \in \omega} \ps^{\leq n}(A)$, basta esibire una successione di enumerazioni $n \mapsto (\alpha_n \twoheadrightarrow \ps^{\leq n}(A))$, per poter usare il corollario sull'unione - cos\`i da evitare AC.\\
	Essendo $|\omega \times A| = \aleph_0$ dall'ipotesi, possiamo fissare  $f : \omega \twoheadrightarrow \omega \times A : x \mapsto (f_1(x),f_2(x))$\footnote{Avendo fissato una $f$ a caso $f_1$ ed $f_2$ sono di fatto funzioni surgettive - entrambe da $\omega$ al rispettivo insieme - a caso.}
	surgettiva (in realt\`a anche bigettiva). A questo punto possiamo costruire la successione di enumerazioni $(\alpha_n)_{n \in \omega}$ - 
	con $\alpha_n$ enumera $\ps^{\leq n}(A)$ per ogni $n \in \omega$ - per ricorsione numerabile prima forma come segue:
	\begin{itemize}
		\item[$\boxed{m = 0}$] In questo caso $\ps^{=0}(A) = \{\emptyset\}$, dunque $\alpha_0$ \`e la funzione che mappa tutti i naturali in 0, cio\`e $\alpha_0 = \{(n,0)\}_{n \in \omega}$.
		\item[$\boxed{m \implies m + 1}$] Data un'enumerazione $\alpha_m$ di $\ps^{\leq m}(A)$, definiamo l'enumerazione $\alpha_{m + 1}$ di $\ps^{\leq m + 1}(A)$ nella maniera seguente:
		\[ 	\alpha_{m + 1} : \omega \rightarrow \ps^{\leq m + 1} : x \mapsto \begin{cases}
			\emptyset &\text{se $x = 0$} \\
			\underbrace{\alpha_m(\overbrace{f_1(x-1)}^{\textcolor{MidnightBlue}{\in \; \omega}})}_{\textcolor{MidnightBlue}{\in \;\ps^{\leq m}(A)}} \cup \{\underbrace{f_2(x - 1)}_{\textcolor{MidnightBlue}{\in A}}\}\;\textcolor{MidnightBlue}{\ne \;\emptyset} &\text{altrimenti}
			\end{cases}
		\]
		\textcolor{MidnightBlue}{di fatto - a parte lasciare un insieme vuoto - stiamo prendendo tutti gli insiemi in $\ps^{\leq m}(A)$ e ci stiamo aggiungendo un elemento di $A$\footnote{Una sorta di shift per la cardinalit\`a.} - va verificato che questa costruzione mantenga la surgettivit\`a e lo faremo a breve.}
	\end{itemize}
	Verifichiamo ora per induzione che la successione che abbiamo costruito ricorsivamente $(\alpha_n)_{n \in \omega}$ di enumerazioni dei $\ps^{\leq n}(A)$ sia effettivamente tale, ovvero che $\alpha_n$ sia una mappa surgettiva da $\omega$ in $\ps^{\leq n}(A)$ per ogni $n \in \omega$.
	\begin{itemize}
		\item[$\boxed{m = 0}$] La successione costante $\alpha_0 = \{(n,0)\}_{n \in \omega}$ \`e banalmente surgettiva.
		\item[$\boxed{m \implies m + 1}$] Per ipotesi induttiva la successione $\alpha_m : \omega \twoheadrightarrow \ps^{\leq m}(A)$ \`e surgettiva, vogliamo verificare che anche $\alpha_{m + 1} : \omega \to \ps^{\leq m + 1}(A)$ lo \`e.
		Dato $Y \in \ps^{\leq s(m)}(A)$ si danno due casi. Se $Y = \emptyset$, allora per costruzione $Y = \alpha_{s(m)}(0) = \emptyset$. Altrimenti siamo nel caso in cui esiste almeno un elemento $y \in Y$.\\
		In questo caso necessariamente $|Y\setminus\{y\}| \leq m$, quindi per ipotesi induttiva - cio\`e $\alpha_m$ surgettiva - $Y \setminus\{y\} = \alpha_m(t)$ per qualche $t \in \omega$.
		Per la surgettivit\`a di $f : \omega \twoheadrightarrow \omega \times A$, la coppia $(t,y)$ \`e immagine di qualche $x \in \omega$, cio\`e $f(x) = (f_1(x),f_2(x)) = (t,y)$.
		Segue quindi che proprio $x+1$ \`e una controimmagine di $Y$:
		\begin{align*}
			\alpha_{m + 1}(x+1) &= \alpha_m(f_1(x)) \cup \{f_2(x)\} &&(\text{definizione $\alpha_{m+1}$})\\
							   &= \alpha_m(t) \cup \{y\} &&(f(x) = (t,y)) \\
							   &= (Y \setminus\{y\}) \cup \{y\} = Y &&(\text{Hp. induttiva})
		\end{align*}
	\end{itemize}
\end{proof}

\paragraph*{Applicazione}\mbox{}\\
Dimostriamo che l'insieme dei numeri reali algebrici $\A_{\RR}$\footnote{Sarebbe $\overline{\QQ} \cap \RR$.} \`e numerabile. Per questa applicazione, assumiamo le propriet\`a elementari di $\RR$.
L'insieme $\A_{\RR}$ \`e definito come l'insieme degli $x \in \RR$ che sono zeri di qualche polinomio a coefficienti razionali:
\[ \A_{\RR} \Mydef \{x \in \RR | \exists p(x) \in \QQ[x]\setminus\{0\} \; p(x) = 0\}
	\]
I numeri reali che non sono algebrici si dicono \vocab{trascendenti} (= $\RR\setminus(\overline{\QQ} \cap \RR)$), siccome - formalmente, vedremo questo risultato in seguito - $\RR$ non \`e numerabile, deduciamo dalla numerabilit\`a di
$\A_{\RR}$ che ci sono numeri reali trascendenti.\\
Dimostriamo, intanto, che l'insieme $\QQ[x]$, dei polinomi a coefficienti razionali nella indeterminata $x$, \`e numerabile.
Possiamo identificare un polinomio:
\[ p(x) = a_0 + a_1x + a_2x^2 + \ldots + a_d x^d
	\]
con l'insieme dei suoi monomi:
\[ p(x) = \{a_0,a_1x,a_2x^2,\ldots,a_dx^d\}
	\]
e ciascun monomio con la coppia (grado, coefficiente):
\[ p(x) = \{(0,a_0),(1,a_1),\ldots,(d,a_d)\}\,\footnote{Pu\`o essere pensata come funzione da $d$ in $\QQ$, o come una funzione da $\omega$ a $\QQ$ a supporto finito.}
	\]
Formalmente, come accade per i numeri, le coppie ordinate, le funzioni, etc., anche i polinomi non sono oggetti atomici della teoria degli insiemi: occorre, in qualche modo, fissare una codifica.
Quella appena descritta \`e una codifica ragionevole. Rappresentando i polinomi in questo modo:
\[ \QQ[x] \subseteq \psf(\omega \times \QQ) \, \footnote{Pi\`u precisamente $\QQ[x] = \bigcup_{n \in \omega} n \times \QQ \subseteq \psf(\omega \times \QQ)$.}
	\]
per cui, essendo che $|\omega \times \QQ| = \aleph_0 \implies |\psf(\omega \times \QQ)| = \aleph_0$, si ha $|\QQ[x]| \leq \aleph_0$.
Inoltre \`e elementare che $\QQ \hookrightarrow \QQ[x]$ (ad esempio $q \mapsto \{(0,q)\}$ \`e una mappa iniettiva che d\`a tutti i polinomi di grado 0), in tal modo si ha anche l'altra disuguaglianza di 
cardinalit\`a e quindi - come al solito per \hyperref[CB]{Cantor-Bernstein} - $|\QQ[x]| = \aleph_0$. Venendo ad $\A_{\RR}$ abbiamo una facile surgezione:
\begin{multline*}
	f : (\QQ[x]\setminus\{0\}) \times \omega \to \A_{\RR} :\\
	 (p,i) \mapsto \text{``la $i$-esima radice di $p$ se questa esiste, altrimenti 0''}
\end{multline*}
Vediamo, per\`o, in maggior dettaglio come si pu\`o rappresentare $f$ mediante una formula insiemistica.
\[
	\text{``$\alpha$ \`e la $i$-esima radice di $p$''} \equiv
	p(\alpha) = 0 \land |\{x \in \RR | x \leq \alpha \land p(x) = 0\}| = |i|
\]
\begin{multline*}
	y = f(p,i) \equiv \text{``$y$ \`e la $i$-esima radice di $p$''} \\
	\land(y = 0 \land \neg \exists \alpha \in \RR \; \text{``$\alpha$ \`e la $i$-esima radice di $p$''})
\end{multline*}
Per separazione esiste, quindi, $f$, e, di conseguenza $|\A_{\RR}| \leq \aleph_0$. La disuguaglianza opposta \`e immediata perch\'e $\QQ \subseteq \A_{\RR}$ - \`e facile scrivere un polinomio in $\QQ[x]$ che abbia come radice un qualsiasi $q \in \QQ$ fissato.

\begin{exercise}[Ogni gruppo finitamente generato \`e al pi\`u numerabile]
	Dato un insieme $X$, una funzione $f : X \times X \to X$, e un sottoinsieme $A \subseteq X$ al pi\`u numerabile, dimostra che esiste un $\ol A \subseteq X$ al pi\`u numerabile tale che $A \subseteq \ol A$\footnote{Typo di Mamino, questa ipotesi manca nelle dispense, ma c'\`e scritta nel 2021.}
	e $f[\ol A \times \ol A] \subseteq \ol A$. Concludi che un gruppo finitamente generato \`e al pi\`u numerabile.
\end{exercise}

\begin{soln}[DA COMPLETARE]
	L'idea per risolvere la prima parte \`e ottenere $\ol A$ come unione numerabile di una successione che parta da $A$ e che ad ogni passaggio nel faccia l'immagine via $f$, in modo tale che l'unione rispetti poi tutte le propriet\`a volute.
	Per questa ragione definiamo per ricorsione numerabile prima forma la seguente successione:
	\[ A_n : \omega \to \ps(X) : \begin{cases}
		A_0 = 0 \\
		A_{n + 1} = f[A_n \times A_n] \cup A_n
	\end{cases}
		\]
	e quindi $\ol A := \bigcup_{n \in \omega} A_n$. Naturalmente, essendo $A_n \subseteq X$ per ogni $n \in \omega$ - per definizione di $f$ -, si ha che $\ol A \subseteq X$; inoltre, presi $x,y \in \ol A$ si ha $x \in A_i$ e $y \in A_j$, per $i, j \in \omega$, poich\'e $A_n \subseteq A_{n + 1}$, cio\`e la successione \`e crescente,
	si ha $x,y \in A_{\max(i,j)}$, per cui per definizione di $f$ si conclude $f(x,y) \in A_{\max(i,j) + 1} \subseteq \ol A$.\\
	Ci rimane da verificare che $|\ol A| \leq \aleph_0$, per fare ci\`o vogliamo usare il corollario dell'unione - che evita AC -, dobbiamo quindi procurarci una 
	successione di enumerazioni $n \mapsto \alpha_n$, con $\alpha_n : \omega \twoheadrightarrow A_n$. Definiamo la successione per ricorsione numerabile.
	\begin{itemize}
		\item[$\boxed{n = 0}$] Per ipotesi $A_0 = A$ \`e al pi\`u numerabile quindi possiamo fissare una sua enumerazione ed usare quella come $\alpha_0$.
		\item[$\boxed{n \implies n + 1}$] Supponiamo di avere per ipotesi induttiva $\alpha_n : \omega \to A_n$ e costruiamo $\alpha_{n + 1} : \omega \to A_{n + 1}$
	\end{itemize}
	Dimostriamo ora per induzione che la successione cos\`i definita \`e effettivamente una successione di enumerazione per gli $A_n$\footnote{Come nella proposizione
	sulle parti finite, non solo la successione di enumerazioni ci permette di applicare il corollario sull'unione, ma ha come conseguenza - cosa che stiamo per dimostrare per induzione - che $\forall n \in \omega \; |A_n| \leq \aleph_0$, che a priori non sapevamo.}.
	\begin{itemize}
		\item[$\boxed{n = 0}$] Abbiamo preso gi\`a $\alpha_0$ enumerazione nelle costruzione per ricorsione, quindi non c'\`e altro da dimostrare.
		\item[$\boxed{n \implies n + 1}$] Supponiamo che $\alpha_n$ sia surgettiva e dimostriamo che lo \`e $\alpha_{n + 1} : \omega \to A_{n + 1}$. Dato $y \in A_{n + 1} = f[A_n \times A_n] \cup A_n$,
		ci sono tre possibilit\`a:
		\begin{itemize}
			\item se $y \in A_n \setminus f[A_n \times A_n]$: in tal caso, essendo $\alpha_n$ surgettiva esiste $x \in \omega$ tale per cui $\alpha_n(x) = y$
			\item se $y \in A_n \cap f[A_n \times A_n]$: in tal caso $y = f(z)$, con $z = (a,b) \in A_n \times A_n$, e di nuovo, per surgettivit\`a di $\alpha_n$ esistono $a',b' \in A_n$, per cui $\alpha_n(a') = a$ e $\alpha_n(b') = b$
			\item se $y \in f[A_n \times A_n] \setminus A_n$:
		\end{itemize}
	\end{itemize}
	Per la seconda parte, dato un gruppo $G$ e $S \subseteq G$, tale che $|S| < \aleph_0$ e $\langle S \rangle = G$, consideriamo la funzione:
	\[ f : G \times G \to G : (x,y) \mapsto xy^{-1}
		\]
	per quanto dimostrato nella prima parte esiste $\overline{S} \subseteq G$ tale che $S \subseteq \ol S$ e $|\ol S| \leq \aleph_0$. Essendo ora $G$ non vuoto e finitamente generato si ha $S \ne \emptyset$,
	per cui $\ol S \ne \emptyset$, verifichiamo dunque che $\ol S$ \`e un sottogruppo di $G$.
	\begin{itemize}
		\item Essendo $\ol S \ne \emptyset$ esiste $g \in \ol S$, e per ipotesi $e_G = gg^{-1} f(g,g) \in \ol S$.
		\item Dal primo punto segue facilmente che $\forall g \in \ol S \; g^{-1} = e_G \cdot g^{-1} = f(e_G,g) \in \ol S$.
		\item Infine $\forall g,h \in \ol S$, per il punto precedente vale sempre che $h^{-1} \in \ol S$, e per ipotesi $gh = f(g,h^{-1}) \in \ol S$.
	\end{itemize}
	Abbiamo quindi che $\ol S$ \`e un sottogruppo di $G$ che contiene $S$, per cui \`e proprio tutto il gruppo, $\ol S = G$, e quindi vale $|G| = |\ol S| \leq \aleph_0$.
\end{soln}

\subsection{Ordini densi numerabili}
Il prossimo risultato che vedremo \`e, come al solito, dovuto a Cantor, e caratterizza l'ordine di $\QQ$ a meno di isomorfismi.

\begin{definition}[Densit\`a]
	Sia $(A,<)$ totalmente ordinato, e $B \subseteq A$, diciamo che $B$ \`e \vocab{denso in} $(A,<)$ se:
	\[ \forall x,y \in A \; x < y \rightarrow \exists z \in B \; x < z < y
		\]
	\textcolor{MidnightBlue}{cio\`e tra due elementi di $A$ c'\`e sempre un elemento di $B$}.\\
	Inoltre diciamo che $(A, <)$ \`e \vocab{denso}, cio\`e \`e denso in se stesso, se:
	\[ \forall x,y \in A \; x < y \rightarrow \exists z \in A \; x < z < y
		\]
	\textcolor{MidnightBlue}{cio\`e tra due elementi di $A$ c'\`e sempre qualche elemento di $A$}.
\end{definition}

\begin{example}[$(\QQ,<)$ \`e denso in se stesso]
	Abbiamo gi\`a osservato, in un esercizio, che $\QQ$ \`e denso, infatti:
	\[ x < y \rightarrow x < \frac{x+y}{2} < y
		\]
	cio\`e presi due qualsiasi elementi di $\QQ$, la loro media aritmetica \`e sempre in mezzo e sta in $\QQ$ (formalmente le due disuguaglianze si giustificano con le operazioni di $\QQ$ + l'ordinamento totale + le propriet\`a 
	di compatibilit\`a tra operazioni e ordinamento).
\end{example}

\begin{notexample}[$(\omega,<)$ non \`e denso in se stesso]
	L'insieme $\omega$ con il suo ordinamento naturale non \`e denso, perch\'e $\not\exists z \in \omega \; 0 < z < 1$.
\end{notexample}

\begin{theorem}[Teorema di isomorfismo di Cantor]
	\label{isoCantor}
	Sia $(A,<)$ un insieme totalmente ordinato tale che:
	\begin{enumerate}[1.]
		\item $|A| = \aleph_0$
		\item $(A,<)$ \`e denso
		\item $(A,<)$ non ha \vocab{estremi}, ossia non ha n\'e massimo n\'e minimo elemento
	\end{enumerate}
	allora $(A,<)\sim(\QQ,<)$.\footnote{Notare che le propriet\`a elencate sono cose che abbiamo gi\`a verificato che valgono per $\QQ$, quindi questo teorema \`e la caratterizzazione di $(\QQ,<)$ come ordine totale.}
\end{theorem}

\textcolor{MidnightBlue}{L'idea \`e di costruire l'isomorfismo per ricorsione. Ad ogni passo della ricorsione avremo $f_i : A_i \rightarrow Q_i$ isomorfismo con $A_i \subseteq A$ finito 
e $Q_i \subseteq \QQ$. Dovremo quindi estendere $f_i$ ingrandendo il suo dominio.
Supponiamo, per esempio, di voler definire $f_{i+1}(x)$ con $x \not \in A_i$. Allora, siccome $A_i$ \`e finito, per sapere la posizione di $x$ a ciascuno degli elemento di $A_i$, ci basta sapere quale sia l'ultimo elemento prima di $x$,
e quale sia il primo dopo $x$ - diciamo che, per esempio, sono $a_2$  e $a_3$ rispettivamente. Dovremo allora mandare $x$ in un $f_{i+1}(x)$ con $f_i(a_2) < f_{i+1}(x) < f_i(a_3)$, e questo esiste per la densit\`a di $\QQ$.}

\begin{figure}[H]
	\centering
	\includegraphics[width = 12cm]{immagini/IsoCantor.png}
\end{figure}

\textcolor{MidnightBlue}{Ragionando simmetricamente, possiamo anche estendere $f_i$, dato un $y \in \QQ$ con $y \not\in Q_i$, in modo tale che $y \in \Imm(f_{i+1})$.}

\begin{figure}[H]
	\centering
	\includegraphics[width = 12cm]{immagini/IsoCantor2.png}
\end{figure}

\textcolor{MidnightBlue}{In definitiva, ci basta quindi fissare un'enumerazione di $A$ e una di $\QQ$, e fare questi passi di estensione in maniera alternata, assicurandoci cos\`i di aggiungere al dominio della $f$, uno per uno,
tutti gli elementi di $A$, e di aggiungere all'immagine, uno per uno, tutti gli elementi di $\QQ$. Ci far\`a comodo la segue osservazione.}

\begin{remark}[L'unione di un insieme di funzioni compatibili \`e una funzione]
	Sia $F \subseteq \ps(A \times B)$ un insieme di funzioni. Se vale che:	
	\[ \forall f_1,f_2 \in F \; f_{1|\Dom(f_1) \cap \Dom(f_2)} = f_{2|\Dom(f_1) \cap \Dom(f_2)}
			\]
	cio\`e se le funzioni da $A$ a $B$ coincidono sull'intersezione dei domini a due a due, $\forall x \in \Dom(f_1) \cap Dom(f_2) \; f_1(x) = f_2(x)$, allora $\bigcup F$ \`e ancora una funzione dall'unione dei domini a $B$:
	\[ \bigcup F : \bigcup\{\Dom(f) | f\in F\} \rightarrow B
		\]
\end{remark}

\begin{proof}
	Bisogna verificare che vale la propriet\`a fondamentale delle funzioni, ovvero che, se $(x,y_1) \in \bigcup F$ e $(x,y_2) \in \bigcup F$, allora $y_1 = y_2$.\\
	Dalla prima cosa abbiamo che esiste $f_1 \in \bigcup F$ tale che $f_1(x) = y_1$ e, dalla seconda, sappiamo che esiste $f_2 \in \bigcup F$ tale che $f_2(x) = y_2$, ma questo significa - per definizione di dominio -
	che $x \in \Dom(f_1) \cap \Dom(f_2)$, dunque dall'ipotesi si ha che:
	\[ y_1 = f_1(x) \overset{\text{Hp.}}{=} f_2(x) = y_2
		\]
\end{proof}

Siamo ora pronti per dimostrare formalmente il teorema.

\begin{proof}
	Per l'ipotesi 1 possiamo fissare un'enumerazione di $A$ e $\QQ$ rispettivamente:
	\[ A = \{a_i | i \in \omega\} \qquad \QQ = \{q_i | i \in \omega\}
		\]
	vogliamo costruire una successione di funzioni $(f_i)_{i \in \omega}$ tale per cui, per ogni $i \in \omega$:
	\begin{enumerate}[(1)]
		\item $f_i : A_i \rightarrow Q_i$ con $|A_i| = |Q_i| < \aleph_0$.
		\item $f_i$ \`e un isomorfismo di ordini fra $A_i$ e $Q_i$.
		\item $f_{i} \subseteq f_{s(i)}$ \textcolor{MidnightBlue}{- ossia $f_{s(i)}$ estende $f_i$}.
		\item $\forall j < i \; a_j \in A_i \land q_j \in Q_i$, \textcolor{MidnightBlue}{ossia il dominio e l'immagine contengono sempre almeno i primi $i$ elementi delle rispettive enumerazioni (poi possono contenere anche altro come vedremo dopo nella costruzione):
		\[ 	\{a_0,\ldots,a_{i-1}\} \subseteq \Dom(f_i) = A_i \qquad Q_i = \{q_0,\ldots,q_{i-1}\} \subseteq \Imm(f_i) = Q_i
			\]}
	\end{enumerate}
	Dando per buona l'esistenza e l'unicit\`a della successione sopra, possiamo definire:
	\[ f := \bigcup_{i \in \omega} f_i
		\]
	\begin{itemize}
		\item[$\diamondsuit$] \underline{$f$ \`e una funzione}: da $(3)$ segue, facilmente per induzione, che $\forall i < j \; f_i \subseteq f_j$, quindi \`e una successione di funzioni ciascuna che estende la precedente, per cui sono compatibili e l'unione \`e ancora una funzione.
		\item[$\diamondsuit$] \underline{$f : A \to \QQ = \Imm(f)$}: da $(4)$ vale che $\Dom(f) = \bigcup_{i \in \omega}\Dom(f_i) = \bigcup_{i \in \omega} A_i = A$ e $\Imm(f) = \bigcup_{i \in \omega}\Imm(f_i) = \bigcup_{i \in \omega} Q_i = \QQ$ - dove abbiamo appunto usato
		che dominio e immagine $i$-esimi, $A_i$ e $Q_i$, per (4) contengono i primi $i$-termini della enumerazione. Poich\'e abbiamo preso come codominio proprio l'immagine $f$ \`e automaticamente surgettiva. 
		\item[$\diamondsuit$] \underline{$f$ strettamente crescente}: dati $x, y \in A$ tali per cui $x < y$, si ha che $x \in A_i$ e $y \in A_j$, per $i,j \in \omega$, ora detta $t:= \max(i,j)$ si ha che $x,y \in A_t$ e, per $(4)$, $A_t \subseteq f(t)$, con $f_t$ isomorfismo - per (2), segue quindi:
		\[ x < y \leftrightarrow f(x) = f_t(x) < f_t(y) = f(y)
			\]
	\end{itemize}
	Non ci resta altro da fare che costruire la successione $(f_i)_{i \in \omega}$ per ricorsione numerabile - prima forma. Poniamo $f_0 = \emptyset$,
	per costruire $f_{s(i)}$ definiamo prima un passo intermedio $f_{i+0.5}$ \textcolor{MidnightBlue}{- nel primo passo, $f_{i + 0.5}$, estendiamo la funzione aggiungendo $a_{i}$ al dominio (se non ci fosse gi\`a), nel secondo passo, $f_{i + 1}$, analogamente, aggiungiamo $q_i$ all'immagine (se non ci fosse gi\`a).}
	Data quindi $f_i : A_i \to Q_i$, distinguiamo due casi:
	\begin{itemize}
		\item \underline{se $a_i \in \Dom(f_i)$}: allora $f_{i + 0.5} = f_i$.
		\item \underline{altrimenti}: detto $\overline{j} := \min\{j \in \omega | f_i \cup \{(a_i,q_j)\} \; \text{\`e un isomorfismo}\}$ \textcolor{MidnightBlue}{(tecnicamente basta strettamente crescente visto che l'immagine \`e il codominio)},
		poniamo $f_{i + 0.5} = f_i \cup \{(a_i,q_{\ol j})\}$.
	\end{itemize}
	Ora possiamo fare il secondo passo di estensione e definire $f_{i+1}$:
	\begin{itemize}
		\item \underline{se $q_i \in \Imm(f_{i + 0.5})$}: allora $f_{i + 1} = f_{i + 0.5}$.
		\item \underline{altrimenti}: detto $\ol \iota := \min\{\iota \in \omega | f_{i + 0.5} \cup (a_{\ol \iota},q_i) \; \text{\`e un isomorfismo}\}$ \textcolor{MidnightBlue}{(come prima basta strettamente crescente)}.\footnote{Notare che la costruzione sarebbe stata simmetrica e funzionante nel caso avessimo voluto costruire $f : \QQ \to A$.}
	\end{itemize}
	Le propriet\`a (1),\ldots,(4) seguono in maniera immediata per induzione, a patto che la costruzione sia ben posta, ossia i minimi esistano. 
	Ad essere precisi, occorre quindi dimostrare, per induzione su $i$, la seguente proposizione:
	\[ \forall i \in \omega \; \text{``la costruzione di $f_i$ \`e ben posta e valgono (1),\ldots,(4)''}
		\]
	Naturalmente il nel caso $f_0 = \emptyset$ la proposizione \`e vera a vuoto, vediamo il passo induttivo, ma solo per la buona definizione.\\
	Per verificare che la definizione di $f_{i + 1}$ sia ben posta, supponendo per ipotesi induttiva che $f_i$ lo sia e che valgano per lei le propriet\`a $(1),\ldots,(4)$,
	bisogna verificare dunque che i minimi esistano in entrambi i passaggi di estensione, e qui entrano in gioco le ipotesi 2 e 3 del teorema. Vediamo che esiste il minimo
	nel primo passaggio di estensione (sar\`a analogo verificarlo nel secondo). Consideriamo:
	\[ \ol j = \min\{j \in \omega | f_i \cup \{(a_i,q_j)\}\,\text{\`e un isomorfismo}\}
		\]
	Per ipotesi induttiva, vale $(1)$ per $f_i$, per cui $A_i$ \`e finito. Detto $n = |A_i|$ (avendo gi\`a fatto il caso $i = 0$ sappiamo che $i > 0$ e dalla costruzione $|i| \leq |A_i|$),
	sfruttando il fatto visto che un ordine totale finito \`e isomorfismo ad un numero naturale possiamo bene ordinare $A_i$ nella maniera seguente:
	\[ A_i = \{\alpha_0,\ldots,\alpha_{n-1}\} \qquad\text{con $\alpha_0<\ldots<\alpha_{n-1}$}
		\]
	Ora, siamo naturalmente nel caso in cui $a_i \not \in \Dom(f_i)$, quindi o $a_i <\alpha_0$, o $\alpha_k < a_i < \alpha_{k+1}$ per
	qualche $k$, o $a_{n-1}<a_i$. Nel primo e terzo caso, rispettivamente, siccome $\QQ$ non ha estremi, c'\`e un $q_j < f_i(\alpha_0)$, o $q_j > f_i(\alpha_n)$, per cui $f_{i + 0.5} \cup \{(a_i,q_j)\}$ \`e un isomorfismo.
	Nel secondo caso, per la densit\`a di $\QQ$, esiste $q_j$ con $f_i(\alpha_k)<q_j<f_i(\alpha_{k+1})$ e quindi di nuovo $f_{i + 0.5}  \cup \{(a_i,q_j)\}$ \`e un isomorfismo. Nella verifica del secondo passo dell'estensione
	faremo lo stesso identico ragionamento con $Q_i$ usando questa volta densit\`a ed illimitatezza di $A$, per trovare gli elementi.
\end{proof}

\begin{corollary}[Ogni ordine al pi\`u numerabile \`e isomorfo ad un sottoinsieme di $\QQ$]
	Sia $(A,<)$ un ordine totale con $|A| \leq \aleph_0$. Allora esiste $B \subseteq \QQ$ tale che $(A,<) \sim (B,<)$ con l'ordinamento indotto su $B$ da $\QQ$.
\end{corollary}

\begin{note}
	Volendo, si potrebbe dimostrare questo corollario ripetendo, con qualche variazione, la dimostrazione del teorema. Ora daremo, per\`o,
	un argomento che, invece, applica il teorema. \`E comodo definire, prima, il prodotto di ordini.
\end{note}

\begin{definition}[Prodotto lessicografico di ordini]
	Dati $(A,<_A)$ e $(B,<_B)$ definiamo il \vocab{prodotto lessicografico} di ordini come:
	\[ (A,<_A) \times (B,<_B) \Mydef (A \times B, <_{A \times B})
 		\]
	dove $(a,b) <_{A \times B} (a',b') \Mydef (b <_B b') \lor (b = b' \land a <_A a')$.
\end{definition}

\textcolor{MidnightBlue}{Ossia $(A,<_A) \times (B,<_B)$ \`e il prodotto cartesiano $A \times B$ munito dell'ordine che confronta prima la seconda componente.}\\
Visualmente, si pu\`o immaginare il prodotto lessicografico $(A,<_A) \times (B,<_B)$ come ``$(A,<_A)$ ripetuto $(B,<_B)$ volte''.

\begin{figure}[H]
	\centering
	\includegraphics[width = 10cm]{immagini/ordine_lessicografico.png}
\end{figure}

\begin{remark}[Il prodotto lessicografico grafico di ordini totali \`e un ordine totale]
	Se $(A, <_A)$ e $(B,<_B)$ sono ordini totali, allora anche $(A,<_A) \times (B,<_B)$ lo \`e.\footnote{\`E una facile verifica che usa appunto la totalit\`a di $<_A$ e $<_B$.}
\end{remark}

\begin{exercise}[Associativit\`a del prodotto lessicografico]
	Dati $(A,<_A),(B,<_B)$ e $(C,<_C)$ dimostra che:
	\[ ((A,<_A) \times (B,<_B)) \times (C,<_C) \sim (A,<_A) \times ((B,<_B) \times (C,<_C))
		\]
	ossia che il prodotto lessicografico di ordini \`e associativo a meno di isomorfismi.
\end{exercise}

Veniamo ora alla dimostrazione del corollario

\begin{proof}
	Se $A = \emptyset$ \`e banale, supponiamo quindi $A \ne \emptyset$ e consideriamo:
	\[ (S,<) \Mydef (\QQ,<) \times (A,<)
		\]
	\begin{figure}[H]
		\centering
		\includegraphics[width = 10cm]{immagini/coroll_iso_cantor.png}
	\end{figure}
	L'insieme $S = \QQ \times A$ \`e numerabile, inoltre, dato $(q,a) \in \QQ \times A$ abbiamo:
	\[ (q-1,a) < (q,a) < (q+1,a) 
		\]
	quindi $\QQ \times A$ non ha estremi - n\'e superiori n\'e inferiori. Per verificare che \`e denso, consideriamo $(q_1,a_1) < (q_2, a_2)$ e verifichiamo che in ogni caso c'\`e sempre un altro elemento di $S$ strettamente nel mezzo.
	\begin{itemize}
		\item Se $a_1 < a_2$, allora $(q_1,a_1) < (q_1+1,a_1) < (q_2,a_2)$ - e funziona per la definizione di ordinamento nel prodotto lessicografico.
		\item Se $a_1 = a_2$ - per cui siamo nella stessa copia di $\QQ$ -, si ha che $(q_1,a_1) < \left(\frac{q_1+q_2}{2},a_1\right) < (q_2,a_2) = (q_2,a_1)$.
	\end{itemize}
	quindi $(S,<)$ \`e denso e per il \hyperref[isoCantor]{teorema di isomorfismo di Cantor} si ha $(S,<) \sim (\QQ,<)$.
	Infine $A \hookrightarrow \QQ \times A : a \mapsto (0,a)$, quindi, componendo l'immersione (strettamente crescente) con l'isomorfismo trovato, abbiamo che $(A,<)$ \`e isomorfo ad un sottoinsieme di $(\QQ,<)$. 
\end{proof}

\begin{exercise}[Isomorfismo di Cantor senza l'ipotesi di illimitatezza]
	Dimostra che se $(A,<)$ \`e denso \textcolor{red}{(ma non necessariamente senza estremi)} e $2 \leq |A| \leq \aleph_0$, allora $(A,<)$ \`e isomorfismo a uno dei seguenti intervalli di $\QQ$:
	\[ [0,1]_{\QQ} \qquad ]0,1]_{\QQ} \qquad [0,1[_{\QQ} \qquad ]0,1[_{\QQ}
		\]
\end{exercise}

\begin{soln}
	Osserviamo che per l'ipotesi di densit\`a $2 \leq |A| \leq \aleph_0 \implies |A| = \aleph_0$\footnote{Stiamo escludendo l'insieme denso con un solo elemento grazie all'ipotesi che $|A| \geq 2$.}. A questo punto, se $A$ \`e senza estremi si ha:
	\[ (A,<) \sim (\QQ,<) \sim ]0,1[_\QQ
		\]
	($]0,1[_\QQ$ \`e totalmente ordinato [ereditariamente dall'ordine di $\QQ$], senza estremi, numerabile [sottoinsieme di $\QQ$ e c'\`e la mappa $n \mapsto \frac 1n$] e denso [basta prendere la media di due elementi e osservare che sta sempre in mezzo per le propriet\`a algebriche di $\QQ$]).\\
	Negli altri casi si osserva che:
	\[ [0,1]_{\QQ} = ]0,1[_{\QQ} \,\cup\, \{0,1\}  \qquad ]0,1]_{\QQ} = ]0,1[_{\QQ} \,\cup\, \{1\} \qquad [0,1[_{\QQ} = ]0,1[_{\QQ} \,\cup\, \{0\}
		\]
	vediamo, ad esempio, nel caso di $A$ con entrambi gli estremi, che, detti questi $a$ e $b$, si ha:
	\[ (A\setminus\{a,b\},<_{|A\setminus\{a,b\}}) \sim (\QQ,<) \sim (]0,1[_{\QQ},<)
		\]
	e analogamente negli altri due casi. Fissiamo un isomorfismo $f$ tra $A \setminus \{a,b\}$ e $]0,1[_\QQ$ ed estendiamolo all'isomorfismo voluto:
	\[ f' : A \to [0,1]_\QQ : x \mapsto \begin{cases}
		0 &\text{se $x = a$} \\
		1 &\text{se $x = b$} \\
		f(x) &\text{altrimenti}
	\end{cases}
		\]
	si verifica facilmente che \`e effettivamente un isomorfismo e negli altri casi si procede in maniera analoga.
\end{soln}

\subsection{Il grafo random}

La tecnica di estendere indefinitamente isomorfismi parziali che ci ha permesso di dimostrare il teorema di isomorfismo di Cantor
si chiama \vocab{back-and-forth}, ed \`e un metodo fondamentale per trovare isomorfismi fra strutture.\\
Cogliamo questa occasione per suggerire un esercizio di applicazione della medesima tecnica che \`e un po' complicato. Si tratta di definire il \vocab{grafo random}
o \vocab{grafo di Rado}.

\begin{definition}[Grafo]
	Un \vocab{grafo} $(V,e)$ sull'insieme di vertici $V$ \`e dato da una relazione $e$ simmetrica \textcolor{MidnightBlue}{($\forall x,y \in V \; (x,y) \in e \leftrightarrow (y,x) \in e$)} e 
	irriflessiva \textcolor{MidnightBlue}{($\forall x \in V \; (x,x) \not\in e$)}.
\end{definition}

\textcolor{MidnightBlue}{L'idea \`e che $V$ pu\`o essere immaginato come un insieme di punti che possono essere connessi da archi. Dire che c'\`e un arco fra $x$ e $y$ equivale a $(x,y) \in e$.\\}
\mbox{}\\
Partiamo da un'idea intuitiva -  chi ha gi\`a seguito un corso di probabilit\`a sapr\`a formalizzare questa cosa in termini precisi. Data una probabilit\`a $p \in \, ]0,1[$ costruiamo un grafo
$G_p$ con insieme di vertici $\omega$ come segue. Per ogni coppia $(i,j) \in \omega \times \omega$ con $i < j$ \textcolor{MidnightBlue}{- cio\`e la seconda componente \`e sempre strettamente pi\`u grande della prima -}
lanciamo una moneta \textbf{che fa testa con probabilit\`a $p$} - tutte queste monete indipendentemente - e, se viene testa, mettiamo un arco fra $i$ e $j$.\\
Potremmo pensare che i grafi $G_{0.01}$ e $G_{0.99}$ debbano venire molto diversi: uno ha l'1\% degli archi possibili, l'altro ha il 99\%, insomma uno \`e quasi vuoto, l'altro quasi completo. \textbf{Avviene,
tuttavia, che, con probabilit\`a 1, questi grafi sono isomorfismi}\footnote{A meno di rinominare i vertici, che \`e quello che diremo nella definizione di isomorfismo.}, dove, per essere precisi dobbiamo definire cosa sia un isomorfismo tra grafi.

\begin{definition}[Isomorfismo fra grafi]
	I grafi $(V_1,e_1)$ e $(V_2,e_2)$ sono \vocab{isomorfi} se esiste una bigezione $f : V_1 \rightarrow V_2$ tale che:
	\[ \forall v,w \in V_1 \; (v,w) \in e_1 \leftrightarrow (f(v),f(w)) \in e_2
		\]
\end{definition}

Vediamo perch\'e. Dati due sottoinsiemi finiti $X$ e $Y$ di $\omega$, e dato un vertice $v \not \in X \cup Y$ la probabilit\`a che $x$ sia connesso da un arco a tutti i vertici di $X$ e a nessuno di quelli di $Y$
\`e $p^{|X|} \cdot (1 - p)^{|Y|}$\footnote{Eventi indipendenti: $v$ \`e connesso ad un vertice di $X$ con probabilit\`a $p$, ed \`e connesso a tutti i vertici di $X$ con probabilit\`a $p^{|X|}$, viceversa non \`e connesso ad alcun vertice di $Y$ con
probabilit\`a $(1-p)^{|Y|}$.} - come che sia, \`e un certo numero $>0$ - e, avendo preso $X$ ed $Y$ finiti, ci sono infiniti $v \in \omega \land v \not \in X \cup Y$.\\
Si capisce, quindi, che con probabilit\`a 1 - ossia certamente - almeno uno di questi $v$ vincer\`a questa lotteria (ne abbiamo infiniti in fondo), ossia sar\`a connesso 
a tutti gli $X$ e a nessuno degli $Y$. Usiamo l'esistenza di questo $v$ per definire un grafo random.

\begin{definition}[Grafo random]
	Il grafo $(\omega,e)$\footnote{Questa definizione \`e per $\omega$, ma la definizione generale di grafo random \`e data dalla stessa propriet\`a per un generico grafo (numerabile) $(G,e)$.} \`e un \vocab{grafo random} se:\footnote{Typo di Mamino, manca ipotesi di finitezza.}
	\[ \forall X,Y \subseteq \omega \land  X \cap Y = \emptyset \land |X|,|Y|<\aleph_0\; \exists v \in \omega \setminus(X \cup Y) \; \underbrace{X \times \{v\} \subseteq e}_{\textcolor{MidnightBlue}{\forall x\in X \; (x,v) \,\in\, e}} \land \underbrace{(Y \times \{v\}) \cap e = \emptyset}_{\textcolor{MidnightBlue}{\not\exists y \in Y \; (y,v) \,\in\, e}}
		\]
	cio\`e se per ogni coppia di sottoinsiemi finiti di vertici disgiunti esiste un vertice fuori dall'unione di questi ultimi, connesso a tutti i vertici di uno ed a nessuno dei vertici dell'altro.
\end{definition}

\begin{exercise}[Esistenza e unicit\`a del grafo random]
	Dimostra che esiste un grafo random, ed \`e unico a meno di isomorfismi - cio\`e indipendentemente da $p \in \, ]0,1[$ viene sempre lo stesso grafo.\footnote{\underline{Hint}: Usare la tecnica del back-and-forth per l'unicit\`a.}
\end{exercise}

\begin{soln}
	Vediamo separatamente esistenza ed unicit\`a.
	\begin{itemize}
		\item[$\boxed{\text{esistenza}}$] Per l'esistenza dobbiamo definire una relazione su $\omega$\footnote{Come per la definizione di grafo random, l'esistenza la si pu\`o dimostrare pi\`u in generale per un qualsiasi insieme numerabile.} che rispetti la propriet\`a che definisce
		un grafo random. Per fare ci\`o ci serviremo per comodit\`a del \href{https://en.wikipedia.org/wiki/BIT_predicate}{\textcolor{purple}{preditcato BIT}} definito come segue:
		\[ \text{BIT}(i,j) = \left\lfloor \frac{i}{2^j}\right\rfloor \mod 2
			\]
		cio\`e BIT$(i,j) = 1$ se la $j$-esima cifra di $i$ espresso in base 2 \`e 1, altrimenti \`e 0. A questo punto possiamo definire la seguente relazione:
		\[ \forall i,j \in \omega \; i \sim y \equiv \text{BIT}(i,j) \lor \underbrace{\text{BIT}(j,i)}_{= \text{BIT}(i,j)^\top}
			\]
		ovvero la simmetrizzazione della relazione BIT. Questa relazione \`e appunto simmetrica e rispetta la propriet\`a del grafo random, dati infatti $X,Y \subseteq \omega$ finiti e disgiunti, consideriamo:
		\[ x := 2^{\max(X,Y) + 1} + \sum_{i \in X} 2^i
			\]
		ci\`o ci assicura che $x \equiv 1 \pmod{2^i}$ per ogni $i \in X$, per cui $\forall i \in X \; (x,i) \in e$, ed al contempo $x \equiv 0 \pmod{2^j}$ per ogni $j \in \omega \setminus (X \cup \{\max(X,Y) + 1\}) \supseteq Y$, cio\`e $(x,j) \not \in e$ per ogni $j \in Y$.
		Naturalmente non vale nemmeno che $(j,x) \in e$ per $j \in Y$ in quanto $x > j$, dunque la cifra $x$-esima di $j$ in base 2 \`e necessariamente 0 ($j < x < 2^x \implies \left\lfloor \frac{j}{2^x}\right\rfloor = 0$).
		\item[$\boxed{\textrm{unicit\`a}}$] Siano $(S,s)$ e $(R,r)$ due grafi random, verifichiamo che sono isomorfi. Fissiamo due enumerazioni di $S$ ed $R$:
		\[ S = \{s_i | i \in \omega\} \qquad R = \{r_i | i \in \omega\}
			\]
		e costruiamo una successione di funzioni $(f_i)_{i \in \omega}$ tale che:
		\begin{enumerate}[(1)]
			\item $f_i : S_i \to R_i$ con $|S_i|,|R_i| < \aleph_0$.
			\item $f_i$ \`e un isomorfismo di grafi.
			\item $\forall i \in \omega \; f_{i} \subseteq f_{i + 1}$.
			\item $\forall i \in \omega \; \{s_0,\ldots,s_{i-1}\} \subseteq \dom(S_i) = S_i \land \{r_0,\ldots,r_{i-1}\} \subseteq \Imm(f_i) = R_i$.
		\end{enumerate}
		Data quindi $(f_i)_{i \in \omega}$, possiamo definire $f := \bigcup_{i \in \omega} f_i$, con $\dom(f) = \bigcup_{i \in \omega} S_i = S$ - il contenimento dal basso c'\`e per (4), quello dall'alto perch\'e come vedremo nella costruzione della successione $\forall i \in \omega \; S_i \subseteq S$.
		Osserviamo ora che:
		\begin{itemize}
			\item[$\diamondsuit$] \underline{$f$ \`e ben definita e surgettiva}: per $(3)$ l'unione delle funzioni $f_i$ \`e una funzione, inoltre, come per il dominio $\bigcup_{i \in \omega} R_i = R$, per cui $f : S \to R$ \`e surgettiva.
			\item[$\diamondsuit$] \underline{$f$ \`e un isomorfismo}: presi $x,y \in S$ vogliamo verificare che valga $(x,y) \in s \leftrightarrow (f(x),f(y)) \in r$. Dato che $x \in S_i$ e $y \in S_j$, per $i,j \in \omega$, preso WLOG $j > i$, abbiamo $x,y \in S_j$, per cui:
			\[ (x,y) \in s \leftrightarrow (f(x),f(y)) = (f_j(x),f_j(y)) \overset{\text{$f_j$ isomorfismo}}{\in} R
				\]
		\end{itemize}
		Non ci resta altro che costruire $(f_i)_{i \in \omega}$ per ricorsione numerabile. Poniamo $f_0 = \emptyset$, ora dato $f_i$ definiamo $f_{i + 1}$ mediante due passi di estensione come segue. Definiamo prima $f_{i + 0.5}$ a partire da $f_i$ nella maniera seguente:
		\begin{itemize}
			\item se $s_i \in \dom(f_i)$, allora $f_{i + 0.5} = f_i$.
			\item altrimenti sia $\ol j := \min\{j \in \omega | f_{i} \cup \{(s_i,r_j)\} \; \text{\`e un isomorfismo}\}$, poniamo $f_{i + 0.5} = f_{i} \cup \{(s_i,r_{\ol j})\}$.
		\end{itemize}
		Ora definiamo $f_{i+1}$ a partire da $f_{i + 0.5}$:
		\begin{itemize}
			\item se $r_i \in \Imm(f_i)$, allora $f_{i + 1} = f_{i + 0.5}$.
			\item altrimenti sia $\ol \iota := \min\{\iota \in \omega | f_{i + 0.5} \cup \{(s_\iota,r_i)\} \; \text{\`e un isomorfismo}\}$, poniamo $f_{i + 1} =f_{i + 0.5} \cup \{(s_{\ol \iota},r_i)\}$.
		\end{itemize}
		A questo punto dimostriamo per induzione che la costruzione sopra sia ben posta in ogni passaggio e che valgano le propriet\`a (1),\ldots,(4). Per $f_0 = \emptyset$ \`e banale, assumiamo ora che $f_i$ sia ben definita e che valgano le propriet\`a (1),\ldots,(4) e verifichiamo che 
		vale la stessa cosa per $f_{i + 1}$. Per la buona definizione dobbiamo verificare che nei passaggi di estensione i minimi esistano sempre, dunque, nel caso in cui $s_i \not \in \dom(f_i)$ verifichiamo che l'insieme 
		$\{j \in \omega | f_{i} \cup \{(s_i,r_j)\} \; \text{\`e un isomorfismo}\}$ \`e non vuoto, per fare ci\`o dobbiamo verificare che esiste $r_j \in R$ tale che $\forall x \in S_i$ se $(x,s_i) \in s$ allora si ha $(f_i(x),r_j) \in r$.
		Essendo $S_i$ finito, l'insieme degli elementi con cui $s_i$ \`e in relazione \`e finito, e quindi a sua volta lo sar\`a l'insieme delle immagini via $f_i$ di tali elementi, chiamiamolo $U$, ora poich\'e $(R,r)$ soddisfa la propriet\`a 
		del grafo random esiste almeno un elemento $u$ in relazione con tutti gli elementi di $U$ e con nessun elemento del vuoto (basta un qualunque sottoinsieme di $R$ finito e disgiunto da $U$), a questo punto ci basta prendere $r_j = u$.\\
		Analogamente, nel secondo passaggio di estensione, nel caso in cui $r_i \not \in \Imm(f_{i+0.5})$, dobbiamo verificare che l'insieme $\{\iota \in \omega | f_{i + 0.5} \cup \{(s_\iota,r_i)\} \; \text{\`e un isomorfismo}\}$ sia non vuoto, ovvero trovare un $s_\iota \in S$ tale che $\forall y \in R$
		per cui $(y,r_i) \in R$, si abbia $(f_{i+0.5}^{-1}(y),s_\iota) \in s$\footnote{Tecnicamente stiamo gi\`a assumendo di aver verificato che $f_{i + 0.5}$ \`e un isomorfismo, ma questo segue banalmente dall'ipotesi induttiva e dalla costruzione.}.
		Come prima, l'insieme degli elementi di $R_{i + 0.5}$ in relazione con $r_i$ \`e finito, e cos\`i sar\`a anche la sua controimmagine via $f_{i + 0.5}^{-1}$, chiamiamola $V$, a questo punto, essendo $(S,s)$ un grafo random,
		esiste un elemento $v \in S$ in relazione con tutti gli elementi di $V$ e nessuno del vuoto, ci basta dunque prendere $s_\iota = v$.\\
		Verifichiamo ora che valgano le propriet\`a (1),\ldots,(4) per $f_{i+1}$. \`E immediato osservare che $f_i \subseteq f_{i + 0.5} \subseteq f_{i + 1}$ per costruzione, e, sempre per costruzione, $s_i \in \dom(f_{i+1})$ e $r_i \in \Imm(f_i + 1)$, in tal modo abbiamo verificato (3) e (4).
		In entrambi i passaggi della costruzione aggiungiamo al pi\`u un elemento, dunque $S_{i+1}$ e $R_{i + 1}$ sono finiti, per cui abbiamo (1), infine, $f_{i+0.5}$ \`e un isomorfismo per costruzione e da questo segue che per costruzione lo \`e anche $f_{i+1}$, dunque abbiamo anche (2).
	\end{itemize}
\end{soln}